\documentclass[UTF8]{ctexart}

\newif\ifBenchTikzExternalize
\BenchTikzExternalizefalse
\InputIfFileExists{bench-config.tex}{}{}

\newcommand{\benchtexttoken}{ALPHA}
\newcommand{\benchfiguretoken}{1.00}
\newcommand{\benchpreambletoken}{P0}

\usepackage{amsmath}
\usepackage{booktabs}
\usepackage{hyperref}
\usepackage{tikz}

\title{W2 CJK Font Workload}
\author{FastXeBench}
\date{}

\begin{document}
\maketitle

\section{背景}
这个工作负载用于模拟 XeLaTeX 在中英文混排场景下的编辑体验。当前标记词为 \benchtexttoken，用于触发正文级别的小修改，并带有一个基础引用路径 \cite{refalpha}。

\section{公式与引用}
我们使用一个简单的目标函数
\begin{equation}
\label{eq:cjk-loss}
\mathcal{L}(x) = \sum_{k=1}^{m} \left(x_k - \mu_k\right)^2 .
\end{equation}
公式~\eqref{eq:cjk-loss} 与后面的图表一起，构成中文论文中常见的最小样式组合。

\section{图形}
\begin{figure}[h]
\centering
\begin{tikzpicture}[scale=\benchfiguretoken]
  \draw[->] (0,0) -- (4.2,0) node[right] {时间};
  \draw[->] (0,0) -- (0,2.8) node[above] {延迟};
  \filldraw[fill=black!15] (0.4,0.4) rectangle (1.1,1.8);
  \filldraw[fill=black!35] (1.5,0.4) rectangle (2.2,2.2);
  \filldraw[fill=black!55] (2.6,0.4) rectangle (3.3,1.4);
  \node[below] at (0.75,0.4) {正文};
  \node[below] at (1.85,0.4) {图形};
  \node[below] at (2.95,0.4) {缓存};
\end{tikzpicture}
\caption{一个带中文标签的简单柱状图。}
\label{fig:cjk-bars}
\end{figure}

如图~\ref{fig:cjk-bars} 所示，中英文混排、数学环境和矢量图可以在同一个最小工作负载中同时出现。

\begin{table}[h]
\centering
\begin{tabular}{@{}lr@{}}
\toprule
场景 & 时间（秒） \\
\midrule
冷启动 & 6.1 \\
稳态编译 & 2.0 \\
\bottomrule
\end{tabular}
\caption{示意性的中文表格。}
\end{table}

\bibliographystyle{plain}
\bibliography{refs}

\end{document}
